\chapter{Metodología}\label{ch:metodologia}

\section{Diseño general}\label{sec:diseno_general}

El objetivo metodológico de esta tesis consiste en caracterizar la distribución espacial y la magnitud de los máximos de precipitación y viento asociados a los ciclones extratropicales del Atlántico Sur, adoptando un marco de referencia lagrangiano centrado en el sistema. 
El análisis se condiciona sistemáticamente por dos factores determinantes: la fase del ciclo de vida y su región de ciclogénesis en Sudamérica. 
La estrategia experimental se estructura en una secuencia lógica: se parte de una base de datos preexistente de trayectorias y una clasificación objetiva de fases evolutivas; sobre esta base, se extraen campos horarios del reanálisis alrededor del centro de cada sistema para construir compositos representativos y, finalmente, se aplican técnicas estadísticas para cuantificar y caracterizar los extremos de las variables de superficie seleccionadas.

\section{Datos del Reanálisis ERA5}\label{sec:datos_era5}

Para la caracterización de los campos atmosféricos se utilizó el reanálisis de quinta generación del ECMWF, ERA5, descrito por \citet{hersbach2020era5}. 
Se analizó el período comprendido entre 1979 y 2020, accediendo específicamente a las variables de precipitación y viento a través del conjunto de datos \textit{``ERA5 hourly data on single levels from 1940 to present''} \citep{hersbach_era5_2023b}. 
La elección de ERA5 de alta resolución resulta crítica para este diseño experimental, ya que permite superar las limitaciones documentadas por \citet{hodges2011comparison}, quienes demostraron que los reanálisis de generaciones previas presentaban discrepancias significativas en la captura de la intensidad máxima de los ciclones en el Hemisferio Sur debido a resoluciones espaciales insuficientes. 
Esta capacidad superior del ERA5 responde directamente a su alta resolución espacial ($\sim$31 km) y temporal (horaria). 
Esta capacidad para resolver escalas finas es crítica en este estudio, ya que, tal como han demostrado \citet{priestley2022improved} y \citet{dalanhese2023new}, la mayor resolución reduce la subestimación de la intensidad de los sistemas observada en generaciones previas de reanálisis, permitiendo una representación más realista de los gradientes de viento y los núcleos de precipitación asociados a los ciclones.

No obstante, debe reconocerse una limitación documentada en la representación de extremos, tal como lo muestra \citet{chen2024evaluation} en la cuantificación del ERA5 y encontraron que subestima sistemáticamente los extremos del top 5\% de ciclones extratropicales, con sesgos del -10.2\% en viento y -22.6\% en precipitación respecto a observaciones satelitales y estaciones superficiales.

\subsection{Variables meteorológicas}\label{subsec:variables}
La caracterización de los patrones espaciales se realizó a partir de campos horarios extraídos de ERA5, seleccionando las variables de superficie:

\begin{itemize}
    \item Precipitación total acumulada.
    \item Magnitud del viento en superficie (10 m).
\end{itemize}

La decisión de analizar los campos de viento y precipitación, y en conjunto, responde a la necesidad de capturar la estructura de los eventos individuales y compuestos (\textit{compound events}). Investigaciones recientes de ciclones extratropicales en latitudes medias del hemisferio norte, como la de \citet{chen2025characteristics}, destacan que la coocurrencia espacio-temporal de extremos de viento y precipitación genera características que difieren significativamente de los extremos aislados. 

En este sentido, según \citet{zscheischler2018future}, la concurrencia varía localmente según la orografía y el patrón meteorológico. Dado que la organización espacial de estos eventos concurrentes es el resultado directo del forzante sinóptico asociado a los ciclones, se propone evaluar ambas variables en conjunto. Esta interacción define el evento más allá de sus efectos individuales y motiva la aplicación de técnicas multivariadas (MEOF) para desacoplar sus modos de variabilidad y analizar la coherencia espacial en los sistemas del Atlántico Sur.

\section{Base de trayectorias y fases}\label{sec:bd_ciclones}

Para la identificación y caracterización de los sistemas, esta investigación utiliza una versión enriquecida de la base de ciclones en el Atlántico Sur. Específicamente, se emplea la base de datos procesada por \citet{coutodesouza2024new}, quien tomó el catálogo original de trayectorias (\textit{tracking}) desarrollado por \citet{gramcianinov2019properties} y le incorporó una clasificación objetiva de las fases evolutivas mediante el algoritmo \textit{CycloPhaser} \citep{deSouza2025CycloPhaser}. De esta forma, el insumo principal de este trabajo combina un seguimiento espacial robusto con un diagnóstico dinámico a lo largo de todo el ciclo de vida del ciclón.

Dado que los fundamentos dinámicos que justifican el uso de la vorticidad relativa y la segmentación del ciclo de vida ya fueron expuestos en las Secciones \ref{sec:tb_vorticidad} y \ref{subsec:tb_fases} de la revision bibliografica, esta sección se limita a describir las especificaciones técnicas y los criterios de filtrado del conjunto de datos utilizado.

En su construcción original \citep{gramcianinov2019properties}, el catálogo base de trayectorias se deriva de los campos del reanálisis ERA5, aprovechando su alta resolución para capturar sistemas de mesoescala y ciclogénesis orográficas, comunes en la región de estudio \citep{gramcianinov2020analysis}. 

Los parámetros de configuración del algoritmo y los filtros aplicados se detallan en la Tabla \ref{tab:tracking_params}.

\begin{table}[ht]
\centering
\caption[Parámetros y filtros de la base de trayectorias]{Parámetros técnicos y filtros aplicados en la base de datos de trayectorias \citep{gramcianinov2020analysis}.}
\label{tab:tracking_params}
\begin{tabular*}{\textwidth}{@{\extracolsep{\fill}}ll}
\hline
Parámetro & Especificación \\ 
\hline
Fuente de datos & ERA5 (ECMWF) \\
Resolución espacial & $\sim$31 km ($0.25^\circ \times 0.25^\circ$) \\
Variable de seguimiento & Vorticidad relativa en 850 hPa ($\zeta_{850}$) \\
% % Nivel de filtrado espectral & T42 (para aislar la escala sinóptica) \\
% % Umbral de intensidad & $\zeta_{850} < -1.0 \times 10^{-5}$ s$^{-1}$ \\
Duración mínima & 24 horas \\
Desplazamiento mínimo & 1000 km \\ 
\hline
\end{tabular*}
\end{table}

% \begin{table}[ht]
% \centering
% \caption{Parámetros técnicos y filtros aplicados en la base de datos de trayectorias \citep{gramcianinov2019properties}.}
% \label{tab:tracking_params}
% \begin{tabular*}{\textwidth}{@{\extracolsep{\fill}}ll}
% \hline
% Parámetro & Especificación \\ 
% \hline
% Fuente de datos & ERA5 (ECMWF) \\
% Resolución espacial & $\sim$31 km ($0.25^\circ \times 0.25^\circ$) \\
% Variable de seguimiento & Vorticidad relativa en 850 hPa ($\zeta_{850}$) \\
% % % Nivel de filtrado espectral & T42 (para aislar la escala sinóptica) \\
% % % Umbral de intensidad & $\zeta_{850} < -1.0 \times 10^{-5}$ s$^{-1}$ \\
% Duración mínima & 24 horas \\
% Desplazamiento mínimo & 1000 km \\ 
% \hline
% \end{tabular*}
% \end{table}


Sobre este seguimiento espacial continuo, la capa de enriquecimiento metodológico de \citet{coutodesouza2024new} le asigna una etiqueta específica a cada paso de tiempo de la trayectoria. Esta clasificación, utilizando \textit{CycloPhaser}, evalúa la tasa de cambio de la vorticidad central, permitiendo aislar con precisión matemática los momentos de desarrollo y decaimiento del sistema. La Tabla \ref{tab:fases_def} resume las cuatro fases evolutivas consideradas en este estudio.

\input{tabelas_es/tab_fases_def.tex}

\section{Estratificación por regiones de ciclogénesis}\label{sec:regiones_ciclogenesis}

Para evaluar si las variables de superficie asociadas a los ciclones podrian varíar según su origen geográfico, la muestra esta estratificada en tres regiones de ciclogénesis. 
La definición espacial de estos dominios sigue estrictamente la taxonomía de densidad propuesta por \citet{gramcianinov2019properties}, delimitando los sectores SBR, LPB y ARG mediante los recuadros de latitud y longitud detallados en la Tabla \ref{tab:regiones_coords}.

\input{tabelas_es/tab_regiones_coords.tex}

Cabe destacar que esta segmentación geográfica se corresponde con los distintos regímenes dinámicos (subtropical/híbrido para SBR, transicional para LPB y baroclínico para ARG) cuyas características físicas y mecanismos de forzamiento fueron discutidos en la Sección \ref{sec:tb_regiones} del Marco Teórico.

% --- DEFINICIÓN DE LAS POBLACIONES DE ESTUDIO ---

\section{Definición de las poblaciones de estudio}\label{sec:poblaciones_estudio}

Para caracterizar la estructura de los ciclones extratropicales y evaluar cuantitativamente el efecto de su intensidad dinámica sobre las variables superficiales, la base de datos de trayectorias fue segmentada en dos niveles jerárquicos: primero se definió una población climatológica base (Población Global) y, a partir de ella, se extrajo el subconjunto p90 basado en el umbral de intensidad de vorticidad relativa en 850 hPa, $\zeta_{850}^{min}$). 
La extracción de estas poblaciones se realizó estratificando los resultados de forma independiente para las tres regiones de ciclogénesis (SBR, LPB y ARG).

\subsection{Población Global: Climatología de referencia}\label{subsec:grupo_global}

El objetivo de la Población Global es establecer el comportamiento morfológico y estadístico medio de los ciclones con un desarrollo típico y completo en la region. 
Para conformarlo, se aplicó un filtro de coherencia evolutiva a la base ciclones, seleccionando únicamente aquellos ciclones que cumplieron con dos criterios fundamentales:

\begin{enumerate}
    \item Ciclo de vida completo: Presencia confirmada secuencial de las cuatro fases de desarrollo: incipiente (\textit{incipient}), intensificación (\textit{intensification}), madurez (\textit{mature}) y decaimiento (\textit{decay}).
    \item Ausencia de fases complejas: Se excluyeron categóricamente los sistemas que desarrollaron fases secundarias, de re-intensificación (e.g., \textit{intensification 2}, \textit{mature 2}) o fases residuales (\textit{residual}).
\end{enumerate}

Este doble filtrado purga el ruido introducido por vórtices efímeros o de evolución errática, buscando garantizar una muestra climatológica homogénea. 
Para cada ciclón de esta Población Global, se identificó el registro temporal correspondiente a su máxima intensidad (el extremo absoluto de vorticidad relativa a 850 hPa, $\zeta_{850}^{min}$) para alinear los eventos y establecer la métrica de intensidad que se utilizará en la segmentación posterior. 

\subsection{Subconjunto de sistemas de alta intensidad (p90)}\label{subsec:subconjuntos_intensidad}

Para evaluar la sensibilidad de los patrones espaciales y la magnitud de las variables frente a la severidad dinámica del sistema, se extrajo de la Población Global un estrato de alta intensidad.
Este grupo comprende el 10\% de ciclones con los valores más negativos de vorticidad relativa mínima en 850 hPa ($\zeta_{850}^{min}$), correspondientes al primer decil de la distribución climatológica. En adelante, este estrato se denominará estrictamente \textbf{subconjunto p90}. La retención exclusiva de este subconjunto de extremos —excluyendo explícitamente los percentiles intermedios o inferiores— responde a una práctica metodológica de compositing ciclónico. \citet{priestley2022improved} aplican explícitamente este mismo umbral estadístico de vorticidad para aislar las estructuras representativas de los ciclones más intensos.
En el contexto específico del Atlántico Sudoccidental, \citet{andrade2024composite} documentan que ciclones intensos (seleccionados usando percentil 10 de presión mínima) exhiben patrones sinópticos diferenciados y de mayores impactos en superficie. 
Esta convergencia metodológica de seleccionar sistemas intensos, permite analizar al subconjunto p90 como una población estadísticamente y físicamente relevante de ciclones extratropicales del Atlántico Sur.

La Tabla \ref{tab:conteo_ciclones} documenta la reducción del tamaño muestral desde la base cruda hasta la conformación de la Población Global y el subconjunto de intensidad p90.

\input{tabelas_es/tab_conteo_ciclones.tex}

\subsection{Estratificación de ciclones}\label{subsec:validacao_selecao}

Un aspecto necesario de esta estratificación es verificar que los sistemas del subconjunto p90 mantengan una evolución temporal coherente. 
El análisis de la fase en la que se registra el pico de vorticidad (Tabla \ref{tab:fase_max_intensidad}) confirma que el subconjunto p90 respeta la estructura dinámica dominante: aproximadamente el 83\% de los casos alcanza su máxima intensidad durante la fase de madurez, frente a menos del 5\% que lo hace durante la fase de decaimiento. 
Este comportamiento de los ciclones más intensos de la base de datos es coherente con el observado en la Población Global y respalda el criterio estadístico del percentil 90 para capturar sistemas sinópticamente maduros y organizados, excluyendo vórtices transitorios o de desarrollo difuso.

\input{tabelas_es/tab_fase_max_intensidad.tex}

\subsection{Estrategia de análisis comparativo}\label{subsec:estrategia_analise}

Sobre el subconjunto p90 de alta intensidad —y comparándolo sistemáticamente con la Población Global de referencia— se aplicará íntegramente el flujo metodológico detallado en las secciones subsiguientes: construcción de campos por fase, extracción de máximos, estimación de distribuciones PDF, análisis PDF espacial (PDFe) de localización y descomposición EOF univariada y multivariada. 
Esta estrategia comparativa entre la climatología completa y el subconjunto de extremos (p90) permitirá evaluar si una mayor intensidad dinámica del vórtice se asocia con cambios significativos en: (i) la magnitud absoluta de los máximos de viento y precipitación, (ii) la distribución espacial relativa de estos máximos, y (iii) los patrones estructurales dominantes. 
Como documentan \citet{andrade2024composite}, los ciclones más intensos presentan estructuras frontales y patrones de impacto distintivos en comparación con los sistemas de menor intensidad, hipótesis que este diseño analítico busca cuantificar para el Atlántico Sur.

\section{Procesamiento Lagrangiano}\label{sec:procesamiento_lagrangiano}

Para caracterizar la estructura de los sistemas independientemente de su ubicación geográfica, se adopta un marco de referencia lagrangiano fundamentado en los principios de seguimiento objetivo de \citet{sinclair1994objective}. 
Este enfoque, basado en el centrado sobre máximos de vorticidad, captura la evolución dinámica de forma más coherente que el seguimiento tradicional por mínimos de presión. 
Como demostró posteriormente \citet{sinclair1997objective}, el uso aislado de la presión central o la vorticidad en un único punto presenta severas limitaciones para inferir la intensidad de los impactos de un ciclón.
Se define un dominio móvil de $20^{\circ} \times 20^{\circ}$ para cada instante $t$ de la trayectoria. Esta técnica de composición permite aislar la señal de mesoescala y extraer los campos horarios de superficie, garantizando la representación física de la extensión espacial de sus brazos frontales.

El análisis se realiza de forma separada para cada combinación de región de ciclogénesis y fase del ciclo de vida, permitiendo estratificar los resultados de manera simultánea por estos dos factores. 
Para reducir la variabilidad de alta frecuencia inherente a los datos horarios y aislar la condición estructural más representativa de cada fase —evitando la contaminación introducida por los instantes de transición entre regímenes dinámicos adyacentes— se aplica un promediado centrado. 
Específicamente, para cada ciclón y fase, se seleccionan los registros horarios correspondientes al núcleo temporal de dicha fase —los dos o tres pasos centrales según su duración— y se promedian sus campos espaciales. 
Este criterio garantiza que el campo resultante refleje la condición arquetípica del régimen en cuestión (incipiente, intensificación, madurez o decaimiento) y no una mezcla de dos estados sucesivos, generando un único campo promedio por ciclón--fase sobre el cual se ejecutarán las técnicas estadísticas subsiguientes.

\section{Caracterización estadística}\label{sec:tecnicas_estadisticas}

\subsection{Extracción de máximos -- densidad de probabilidad (PDF)}\label{subsec:extraccion_pdf}

A partir del campo promedio por ciclón--fase, se delimita el área de análisis aplicando una máscara circular de $9.5^{\circ}$ de radio, centrada en el sistema como se explica en \ref{sec:procesamiento_lagrangiano}. 
Para la extracción de los máximos, y con el objetivo de mitigar el impacto del ruido o anomalías aisladas a escala de punto de grilla, se aplicó un suavizado espacial bidimensional sobre los campos utilizando un filtro gaussiano ($\sigma = 0.25$) mediante la biblioteca \texttt{scipy.ndimage} de Python. 
Dentro de este radio de influencia continuo y suavizado, se identifica la magnitud del valor máximo absoluto de la variable de interés ($x_i$), ignorando valores fuera del límite espacial, junto con su ubicación en coordenadas relativas al centro, denotada como $\mathbf{s}_i=(\Delta x_i,\Delta y_i)$. 
Este procedimiento genera, para cada combinación de región, fase evolutiva y variable, un conjunto de $m$ valores máximos empíricos y sus respectivas posiciones espaciales (un par magnitud-ubicación por cada evento).

Para caracterizar la distribución estadística de las magnitudes máximas extraídas ($\{x_i\}_{i=1}^{m}$), se estimó la función de densidad de probabilidad (PDF) mediante la técnica de estimación de densidad por kernel (KDE). 
Este procedimiento analítico fue implementado utilizando la clase \texttt{gaussian\_kde} de la biblioteca \texttt{scipy.stats}. 

La función estimada $\hat{f}(x)$ se define como:

\begin{equation}
    \hat{f}(x) = \frac{1}{m h} \sum_{i=1}^{m} K\left(\frac{x - x_i}{h}\right),
\label{eq:kde_pdf}
\end{equation}

donde $h$ es el parámetro de suavizado o ancho de banda (\textit{bandwidth}), y $K(\cdot)$ representa la función núcleo, y se seleccionó un núcleo gaussiano, expresado como:

\begin{equation}
    K(u) = \frac{1}{\sqrt{2\pi}} \exp\left(-\frac{u^2}{2}\right),
\label{eq:gaussian_kernel}
\end{equation}

donde el argumento $u = (x - x_i)/h$ representa la distancia estandarizada entre el valor a evaluar ($x$) y la observación empírica ($x_i$). 
Conceptualmente, la función núcleo ($K$) actúa como un factor de ponderación que toma cada observación individual y distribuye su contribución de probabilidad sobre su entorno inmediato adoptando forma de campana. 
El parámetro $h$ controla directamente la amplitud de esta distribución local; en este trabajo, la biblioteca de cálculo determinó el valor óptimo de $h$ de forma automática aplicando la regla de estimación de Scott. 

La adopción de PDF permite una representación geométrica continua y más analítica de los datos, superando las limitaciones de los histogramas tradicionales. 
Como señala \citet{weglarczyk2018kernel}, la estimación por núcleos evita el sesgo inherente a la selección del número, ancho y punto de inicio de las clases (\textit{binning}), ya que preserva y utiliza la ubicación exacta de cada observación individual. 
Este análisis probabilístico permite contrastar la distribucion de las curvas para evaluar cuantitativamente cómo evolucionan las magnitudes de las variables a lo largo del ciclo de vida del ciclón y cómo difieren según la severidad del sistema y su región de génesis.

\subsection{Distribución espacial: PDF bidimensional}\label{subsec:kde_localizacion}

Para caracterizar la distribución espacial preferencial de la ocurrencia de estos máximos, se utiliza el conjunto de $m$ ubicaciones relativas $\{\mathbf{s}_i\}_{i=1}^{m}$ extraídas en el paso inicial.
A partir de estas posiciones, se extiende la formulación probabilística unidimensional definida previamente (Ecuación \ref{eq:kde_pdf}) para estimar un campo continuo de densidad espacial mediante un PDF bidimensional (PDFe), calculado como el promedio de núcleos centrados en cada máximo:

\begin{equation}
\hat f(\mathbf{s})=\frac{1}{m}\sum_{i=1}^{m} K_h(\mathbf{s}-\mathbf{s}_i),
\label{eq:kde_simple}
\end{equation}

donde $\mathbf{s}=(\Delta x,\Delta y)$ es una posición genérica dentro del dominio centrado y $h$ es el parámetro de suavizado espacial. En este trabajo tambien se adopta un núcleo gaussiano bidimensional:

\begin{equation}
K_h(\mathbf{u})=\frac{1}{2\pi h^2}\exp\!\left(-\frac{\|\mathbf{u}\|^2}{2h^2}\right).
\label{eq:gauss_kernel_2d_simple}
\end{equation}

El resultado es un campo continuo de densidad de probabilidad relativa que indica, para cada posición dentro del dominio centrado, la tendencia espacial a encontrar el máximo de la variable. 
Desde el punto de vista metodológico, la implementación de un PDF bidimensional presenta ventajas sustanciales frente a los métodos discretos de conteo; de acuerdo con \citet{weglarczyk2018kernel}, este enfoque elimina por completo la necesidad subjetiva de definir el tamaño y la orientación geométrica de las celdas. Esta técnica de suavizado continuo resulta idónea para mapear de forma objetiva asimetrías en la distribución espacial de los impactos. En este sentido, \citet{gramcianinov2023impact} demostraron la relevancia de analizar la ubicación relativa de los extremos en un marco de referencia centrado en el sistema, evidenciando que los impactos más severos tienden a confinarse en sectores preferenciales (como el oeste y noroeste en el Atlántico Sur) y que esta configuración espacial asimétrica varía fuertemente con la fase del ciclo de vida del ciclón.

\subsection{Funciones ortogonales empíricas (EOF) univariado}\label{subsec:eof_univariado}

Para identificar los patrones espaciales dominantes de variabilidad en la estructura de los ciclones dentro del marco centrado, se aplica un análisis de Funciones Ortogonales Empíricas (EOF). 
Aunque tradicionalmente esta técnica se aplica a campos de variación suave, investigaciones recientes como la de \citet{sun2022evaluation} han demostrado la viabilidad y robustez de modelos basados en funciones ortogonales para extraer patrones espaciales y su evolución temporal en datos de precipitación (en Corea del Sur).
Este análisis se realiza de manera separada para cada variable (V10, precipitación), región de ciclogénesis y fase del ciclo de vida, utilizando el conjunto de $m$ campos promedio por ciclón--fase.

Previo al análisis, a cada campo se le resta la media de la Población Global del ensamble en cada punto de grilla para obtener las anomalías estructurales. Matemáticamente, si $X(\tau, \mathbf{s})$ representa el campo de la variable para el ciclón $\tau$ (con $\tau = 1, 2, \dots, m$) en la posición espacial relativa $\mathbf{s} = (x, y)$, el campo de anomalías $X'(\tau, \mathbf{s})$ se define como:
\begin{equation}
    X'(\tau, \mathbf{s}) = X(\tau, \mathbf{s}) - \frac{1}{m}\sum_{\tau=1}^{m}X(\tau, \mathbf{s})
\end{equation}
A partir de esta matriz de anomalías se calcula la matriz de covarianza espacial, cuya diagonalización permite descomponer el campo original en una suma lineal de $K$ modos ortogonales espaciales ($EOF_k$) y sus correspondientes amplitudes o componentes principales ($PC_k$):
\begin{equation}
    X'(\tau, \mathbf{s}) = \sum_{k=1}^{K} PC_k(\tau) \cdot EOF_k(\mathbf{s})
\end{equation}

Cabe destacar que no se aplica una remoción de tendencia temporal (\textit{detrend}) a la secuencia de eventos. Dado que el eje muestral está constituido por un índice de eventos discretos separados por intervalos de tiempo irregulares a lo largo del período 1979--2020, la aplicación de un ajuste lineal cronológico estándar asumiría erróneamente un espaciado temporal constante. Adicionalmente, desde un punto de vista físico, el objetivo metodológico es aislar las características morfológicas de la variabilidad inter-sistémica asociadas a las distintas fases del ciclón, por lo que aislar o eliminar las tendencias climáticas de muy baja frecuencia escapa al diseño fundamental de este análisis.

Este método, adaptado al dominio lagrangiano, permite aislar y comparar las estructuras espaciales más recurrentes o típicas asociadas a los ciclones en cada contexto (región/fase), siguiendo la lógica de análisis de patrones aplicada en estudios sinópticos \citep{vera2002cold}.

\subsection{Significancia espacial: Bootstrap-t}\label{subsec:bootstrap_significancia}

Para aislar las señales físicas del ruido estocástico y evaluar la robustez de los patrones espaciales (autovectores) derivados del análisis EOF, se implementó el método de remuestreo no paramétrico \textit{Bootstrap-t}, desarrollado originalmente por \citet{efron1993introduction}. Siguiendo la metodologia expuesta por \citet{hesterberg2015bootstrap} su correcta aplicación en este enfoque resulta óptimo para corregir dinámicamente la asimetría inherente a las distribuciones de precipitación y viento. Se generaron $R = 15,000$ submuestras con reemplazo, garantiza que el error de Monte Carlo sea insignificante al estimar las colas de la distribución.

Para cada iteración $b$, se recalculó el análisis EOF y se obtuvo el estadístico casi pivotal $t^*_b$:
\begin{equation}
    t^*_b = \frac{\hat{\theta}^*_b - \hat{\theta}}{SE^*_b}
\end{equation}
donde $\hat{\theta}$ representa la carga espacial original, $\hat{\theta}^*_b$ su estimación en la submuestra y $SE^*_b$ su error estándar estimado.

Los intervalos de confianza al 95\% se calcularon invirtiendo los cuantiles empíricos $q_{\alpha/2}$ y $q_{1-\alpha/2}$ (con $\alpha = 0.05$) de la distribución de $t^*$:
\begin{equation}
    (\hat{\theta} - q_{1-\alpha/2}SE, \hat{\theta} - q_{\alpha/2}SE)
\end{equation}
Los puntos de grilla se consideraron estadísticamente significativos si su intervalo de confianza excluía el cero. 
Todo el procesamiento estadístico y la descomposición matricial se ejecutaron en el entorno computacional Python. 
Para la gestión de los datos multidimensionales se empleó la biblioteca \texttt{xarray} \citep{hoyer2017xarray}, mientras que la extracción de los modos espaciales y la evaluación de su significancia mediante \textit{bootstrapping} se implementaron utilizando \texttt{xeofs} \citep{rieger2024xeofs}.

\subsection{Análisis de las Componentes Principales}\label{subsec:analisis_pc}

Además de la descomposición espacial tradicional de los campos mediante Funciones Ortogonales Empíricas (EOF), se realizó un análisis de las propiedades estadísticas de las Componentes Principales (PC) asociadas a cada modo. 
Este análisis permite caracterizar la variabilidad y conexiones físicas entre la estructura espacial en superficie y la dinámica de los ciclones.
La PC1, que representa la amplitud temporal del modo dominante de variabilidad espacial, fue caracterizada mediante métricas descriptivas de su distribución de probabilidad. 
Dado que la PC1 captura la proyección de cada ciclón sobre el patrón estructural más recurrente de la región, su análisis estadístico permite inferir propiedades de la Población Global que no son evidentes desde el patrón espacial.

Se calcularon los siguientes estadísticos para la serie de PC1 de cada variable, región y fase:
La asimetría (skewness, $\gamma_1$) se calculó como:

\begin{equation}
    \gamma_1 = \frac{1}{m} \sum_{\tau=1}^{m} \left(\frac{PC_{1,\tau} - \overline{PC_1}}{\sigma_{PC_1}}\right)^3
\end{equation}

donde $m$ es el tamaño muestral, $\overline{PC_1}$ es la media de la serie y $\sigma_{PC_1}$ su desviación estándar. La asimetría positiva indica una cola derecha extendida, es decir, una subpoblación de ciclones con amplitud anormalmente alta en el modo dominante; asimetría negativa indica el patrón inverso.

La curtosis ($\gamma_2$) se obtuvo mediante:

\begin{equation}
    \gamma_2 = \frac{1}{m} \sum_{\tau=1}^{m} \left(\frac{PC_{1,\tau} - \overline{PC_1}}{\sigma_{PC_1}}\right)^4 - 3
\end{equation}

valores positivos (leptocurtosis) indican distribuciones con colas pesadas y concentración central pronunciada, sugiriendo que los ciclones se ajustan estrechamente al patrón arquetípico de una subpoblación de sistemas con estructuras intensas o atípicas respecto al patrón dominante; valores negativos (platicurtosis) indican que los eventos extremos son menos frecuentes y la población de ciclones presenta mayor homogeneidad estructural, con la mayoría de los sistemas agrupándose cerca del comportamiento medio.

\subsection{Relación entre PC1 e intensidad dinámica del ciclón}\label{subsubsec:pc1_vorticidad}

Para establecer la conexión física entre la estructura de los campos de superficie y la intensidad sinóptica del sistema, se evaluó la correlación entre la PC1 de cada variable (precipitación y viento) y el mínimo de vorticidad relativa a 850 hPa alcanzado por cada ciclón durante su ciclo de vida ($\zeta_{850}^{min}$).

La correlación de Pearson se calculó como:

\begin{equation}
    \rho(PC_1, \zeta_{850}^{min}) = \frac{\text{cov}(PC_1, \zeta_{850}^{min})}{\sigma_{PC_1} \cdot \sigma_{\zeta}}
\end{equation}

Esta relación permite responder si el modo dominante de variabilidad espacial se intensifica sistemáticamente con la profundidad del vórtice ciclónico, o si por el contrario la estructura de precipitación y viento responde a forzantes adicionales independientes de la intensidad del centro del ciclón. Una correlación significativa y negativa (dado que $\zeta_{850}^{min}$ es negativa en el Hemisferio Sur) indicaría que los ciclones más intensos presentan mayor proyección sobre el patrón estructural dominante; una correlación débil sugeriría que la organización espacial de los impactos está modulada por factores locales (orografía, inestabilidad convectiva, interacción con frentes preexistentes) y que no necesariamente escalan linealmente con la vorticidad del núcleo.

\subsection{Coherencia entre modos dominantes de precipitación y viento}\label{subsubsec:coherencia_pc1}

Para evaluar si los patrones espaciales dominantes de precipitación y viento varían conjuntamente o responden a dinámicas independientes, se calculó la correlación entre las PC1 de ambas variables dentro de cada región y fase evolutiva:

\begin{equation}
    \rho = \text{corr}(PC_1^{\text{precip}}, PC_1^{\text{viento}})
\end{equation}

Esta correlación cuantifica el grado de acoplamiento entre la variabilidad estructural de ambos campos. Valores cercanos a +1 indican que los extremos de precipitación y viento se organizan espacialmente de manera coherente (eventos compuestos bien estructurados); valores cercanos a 0 sugieren que la variabilidad espacial de precipitación está controlada por otros procesos locales (convección, orografía) independientes de la configuración del campo de viento a escala sinóptica.

% Adicionalmente, se construyeron diagramas de dispersión de $PC_1^{\text{precip}}$ vs $PC_1^{\text{viento}}$ para identificar cuadrantes predominantes y detectar la existencia de ciclones con comportamientos atípicos (alta PC en una variable, baja en la otra), que podrían indicar desacoplamiento entre los procesos dinámicos y termodinámicos del sistema.

\subsection{EOF combinado (multivariado) precipitación--viento}\label{subsec:eof_multivariado}

Para evaluar la covariabilidad espacial entre los campos de precipitación y viento, se realiza un análisis EOF combinado (MEOF). Dado que estas variables poseen unidades físicas y rangos de varianza dispares (mm/h vs. m/s), es un requisito metodológico fue estandarizar las series antes de su integración para evitar que la variable con mayor magnitud numérica domine los modos resultantes \citep{wilks2019statistical_ch13}. 

El procedimiento analítico consta de tres pasos fundamentales:

(i) Calcular las anomalías de cada campo respecto a su media de la Población Global.

(ii) Normalizar dichas anomalías dividiéndolas por la desviación estándar espacialmente promediada de cada variable, asegurando un peso equiparable en la matriz de covarianza combinada. 
Esto con el fin de preservar los gradientes espaciales nativos de cada campo físico, este factor se define como un escalar global único. 
Matemáticamente, el campo estandarizado $Z_{V}(\tau, \mathbf{s})$ para la variable $V$ en el ciclón $\tau$ se obtiene como:

\begin{equation}
    Z_{V}(\tau, \mathbf{s}) = \frac{V'(\tau, \mathbf{s})}{\frac{1}{A}\int_{A} \sigma_{V}(\mathbf{s}) \, d\mathbf{s}}
\end{equation}

donde $V'$ es el campo de anomalías, $\sigma_{V}(\mathbf{s})$ es la desviación estándar temporal en el punto $\mathbf{s}$, y $A$ es el área total del dominio.

(iii) Concatenar espacialmente los campos estandarizados para formar un vector de estado combinado o ampliado, siguiendo el esquema de análisis de patrones acoplados descrito por \citet{bretherton1992intercomparison}:

\begin{equation}\label{eq:meof_vector}
    \mathbf{Y}(\tau) = \big[ Z_P(\tau), Z_W(\tau) \big]
\end{equation}

El análisis EOF aplicado a esta matriz combinada $\mathbf{Y}(\tau)$ permite obtener los modos acoplados que explican la varianza conjunta del sistema, siguiendo la base matemática de descomposición multivariada para eventos compuestos validada por \citet{liang2018multivariate}. La aplicación de esta técnica es instrumental para responder a interrogantes recientes planteadas por \citet{sinclair2023relationship}, permitiendo revelar si los núcleos de precipitación intensa se co-ubican sistemáticamente con los gradientes de viento más fuertes y cómo esta coherencia estructural evoluciona a lo largo del ciclo de vida.

\subsection{Prototipos conceptuales}\label{subsec:prototipos}

Para cumplir con el Objetivo Específico 4, se desarrolló un procedimiento de síntesis visual que integra, en una representación compuesta única por región y fase evolutiva, cuatro fuentes de información estadística derivadas de los análisis previos: la distribución espacial de probabilidad de los máximos de cada variable (PDFe, Sección~\ref{subsec:kde_localizacion}) y el patrón espacial dominante de variabilidad estructural conjunta de las variables (EOF multivariado, Sección~\ref{subsec:eof_multivariado}). 
En este análisis, el procedimiento se realiza para cada combinación región--fase, tanto para la Población Global y la muestra del subconjunto de ciclones intensos (p90), preservando así la representatividad de los patrones obtenidos.

\subsubsection{Superposición de capas estadísticas}

Cada prototipo está compuesto por cuatro capas estadísticas superpuestas en un único panel de coordenadas lagrangianas $(\Delta x, \Delta y)$, centrado en el núcleo del ciclón dentro del dominio de $9.5^{\circ}$ de radio definido en la Sección~\ref{subsec:extraccion_pdf}.
Las cuatro capas son:

\paragraph{Capa 1 y 2 — Isolíneas del PDFe de ubicación de máximos (precipitación y viento).}

A partir de la superficie de densidad $\hat{f}(\mathbf{s})$ calculada según la Ecuación~\ref{eq:kde_simple} para cada variable, se extrajeron dos niveles de contorno correspondientes a los percentiles 75 y 90 de la distribución acumulada de la densidad. 
El procedimiento de cálculo de estos niveles consiste en ordenar todos los valores de la superficie de densidad de menor a mayor, calcular la probabilidad acumulada normalizada sobre ese vector ordenado, e identificar el valor de densidad $\hat{f}^*$ que corresponde a la fracción de probabilidad deseada. Formalmente:

\begin{equation}
    \hat{f}^*_q = \hat{f}_{(k)}, \quad \text{donde } k = \min\left\{j \; : \; 
    \frac{\sum_{i=1}^{j}\hat{f}_{(i)}}{\sum_{i=1}^{m}\hat{f}_{(i)}} \geq q \right\},
\label{eq:nivel_contorno_kde}
\end{equation}

siendo $\hat{f}_{(i)}$ los valores de la densidad ordenados ascendentemente y $q$ la fracción acumulada deseada ($q = 0.75$ o $q = 0.90$). 
La isolínea correspondiente al percentil 75 delimita la región que concentra el 75\% de la masa de probabilidad acumulada de la densidad, y por tanto encierra un área más extensa; la del percentil 90 
encierra el núcleo más denso de ocurrencia de máximos, representando la zona de mayor concentración espacial de los valores extremos. 

Adicionalmente, sobre cada superficie PDFe se identifica la coordenada modal $(\Delta x_{\text{moda}}, \Delta y_{\text{moda}})$, definida como la posición del pico máximo de densidad:

\begin{equation}
    (\Delta x_{\text{moda}}, \Delta y_{\text{moda}}) = 
    \underset{(\Delta x, \Delta y)}{\arg\max} \; \hat{f}(\Delta x, \Delta y),
\label{eq:moda_kde}
\end{equation}

y se señaliza en el panel mediante un marcador estelar. 
Esta coordenada representa la ubicación relativa al centro del ciclón donde estadísticamente es más frecuente encontrar el máximo de la variable analizada.

\paragraph{Capa 3 y 4 — Isolíneas del loading espacial del MEOF$_1$ (precipitación y viento).}

Sobre el mismo panel se superponen las isolíneas de los \textit{loadings} espaciales por variable del primer modo del análisis EOF multivariado (MEOF$_1$) descrito en la Sección~\ref{subsec:eof_multivariado}. 
Dado que el vector de estado combinado $\mathbf{Y}(\tau) = [Z_P(\tau), Z_W(\tau)]$ concatena ambas variables en un único autovector (Ecuación~\ref{eq:meof_vector}), el MEOF$_1$ produce simultáneamente un 
\textit{loading} espacial para precipitación ($\text{EOF}_{1,P}$) y otro para viento ($\text{EOF}_{1,W}$), ambos extraídos del mismo modo dominante de covariabilidad conjunta. 
Son estos dos sub-patrones acoplados los que se grafican en el prototipo, preservando la coherencia espacial entre ambas variables impuesta por la descomposición multivariada y garantizando que las estructuras visualizadas reflejen el modo de variabilidad que maximiza simultáneamente la varianza conjunta de precipitación y viento.

A partir del campo de cada \textit{loading} espacial ($\text{EOF}_{1,P}$ y $\text{EOF}_{1,W}$), se extraen dos conjuntos de niveles de contorno mediante percentiles aplicados separadamente sobre los valores positivos y negativos del autovector: los percentiles 33 y 66 de los valores positivos (representados con línea sólida) y los percentiles 33 y 66 del valor absoluto de los valores negativos (representados con línea discontinua), aplicando previamente un suavizado gaussiano ($\sigma = 2$). Esta estrategia 
de umbralización relativa garantiza que los contornos reflejen la estructura interna del patrón dominante de covariabilidad independientemente de la magnitud absoluta del \textit{loading}, haciendo comparables los patrones entre distintas combinaciones región--fase. Las regiones de carga positiva indican sectores donde la variable tiende a presentar anomalías positivas respecto a la media climatológica del ensamble cuando el 
modo dominante se activa; las regiones de carga negativa indican supresión relativa de la variable en ese mismo contexto dinámico.

\subsubsection{Estructura de la figura de prototipos}

El resultado es una figura matricial de $3 \times 4$ paneles (3 regiones de ciclogénesis $\times$ 4 fases evolutivas) donde cada panel integra simultáneamente las cuatro capas descritas: PDFe de precipitación (contornos en azul), PDFe de viento (contornos en rojo), $\text{EOF}_{1,P}$ de precipitación (isolíneas en verde) y $\text{EOF}_{1,W}$ de viento (isolíneas en naranja). 
La superposición de estas cuatro capas en el espacio lagrangiano centrado en el ciclón permite identificar visualmente el grado de acoplamiento espacial entre la variabilidad estructural dominante multivariada (MEOF) y la distribución de probabilidad de los máximos (PDFe): una alta coincidencia espacial entre el núcleo del PDFe y la región de carga positiva del EOF indica que el patrón dominante de variabilidad está directamente asociado a la ocurrencia de los valores extremos, mientras que un desacoplamiento entre ambas capas señala la existencia de procesos locales o asimetrías frontales que organizan los impactos en 
sectores distintos al patrón modal. 
La coordenada de cada variable, anotada en la leyenda interna de cada panel como $(\Delta x, \Delta y)$ en grados, cuantifica explícitamente el desplazamiento preferencial de los máximos respecto al centro del vórtice ciclónico. Esta síntesis constituye la representación más compacta y completa de la huella de severidad de los ciclones extratropicales del Atlántico Sur, integrando en un único diagrama las dimensiones espacial, probabilística y modal de los campos de precipitación y viento en superficie.
